
\textbf{Proposition 6.2.2.} \uline{Pour toute section positive} @@M0@@ \uline{de} @@M1@@\uline{, le faisceau étale en groupes} @@M2@@ \uline{des commutateurs de} @@M3@@ \uline{est représentable par}
@@M4@@.

\textbf{Corollaire 6.2.3.} \uline{Toute représentation}
@@M5@@
\uline{est de la forme}
@@M6@@
\uline{où} @@M7@@, @@M8@@ \uline{et} @@M9@@ \uline{sont des sections de} @@M10@@ \uline{uniquement déterminés par} @@M11@@.

\textbf{Corollaire 6.2.6.} \uline{Avec les notations de} 6.1. \uline{on a}
@@M12@@
\uline{où} @@M13@@ \uline{sont des sections de} @@M14@@ \uline{uniquement déterminées par} @@M15@@.

\textbf{Remarque 6.2.7.} On peut déterminer les exposants @@M16@@ dans (6.2.5) et (6.2.6) en substituant successivement dans (6.2.5)
@@M17@@, @@M18@@, pour @@M19@@ @@M20@@, d'où
@@M21@@
Par exemple on déduit ainsi de (6.2.6)
@@M22@@
et
@@M23@@
etc.

